\documentclass[11pt]{article}
\usepackage{mathblatt}
% gesetzt von werkzeuge/setzer.py; Lösungen nach bau/bauregeln.md „Lösungen“.
\newcommand{\kennung}{M74}
\begin{document}
\pfheftstil
\fancyfoot[C]{{\footnotesize\color{mbgrau}\kennung\hfill\thepage}}
\pfheftkopf{Hypotenuse berechnen}{Klasse 9 \textperiodcentered{} Lösungen}

\begin{pfloesung}{}
\lz{1a)}{$c$}{gegenüber dem rechten Winkel}{}
\lz{b)}{$r$}{rechter Winkel unten rechts}{}
\lz{c)}{$s$}{rechter Winkel oben, Hypotenuse waagerecht}{}
\lz{d)}{$e$}{rechter Winkel unten, Hypotenuse oben}{}
\end{pfloesung}
\begin{pfloesung}{}
\lz{2a)}{$25$ Kästchen}{$9 + 16 = 25$}{}
\lz{b)}{$25$ Kästchen}{$5 \cdot 5 = 25$}{}
\lz{c)}{nein}{$3 + 4 = 7 \ne 5$; die Flächen passen, die Längen nicht}{}
\end{pfloesung}
\begin{pfloesung}{}
\lz{3a)}{$w^2 = u^2 + v^2$}{Hypotenuse $w$}{}
\lz{b)}{$x^2 = y^2 + z^2$}{Hypotenuse $x$}{}
\lz{c)}{$\overline{BC}^2 = \overline{AB}^2 + \overline{AC}^2$}{rechter Winkel bei $A$, Hypotenuse $\overline{BC}$}{}
\lz{d)}{gilt nicht}{kein rechter Winkel}{}
\end{pfloesung}
\begin{pfloesung}{}
\lz{4b)}{$13$ cm}{$c^2 = 5^2 + 12^2 = 169$ \ $\Rightarrow$ \ $c = 13$; \ $12 < 13 < 17$}{}
\lz{c)}{$15$ cm}{$c^2 = 9^2 + 12^2 = 225$ \ $\Rightarrow$ \ $c = 15$}{}
\end{pfloesung}
\begin{pfloesung}{}
\lz{5.}{$t = \sqrt{r^2 + s^2}$}{$t$ liegt dem rechten Winkel gegenüber: $t^2 = r^2 + s^2$ \ $\Rightarrow$ \ Wurzel}{}
\end{pfloesung}
\begin{pfloesung}{}
\lz{6b)}{$\sqrt{117} \approx 10{,}8$ cm}{$6^2 + 9^2 = 117$ \ $\Rightarrow$ \ $\sqrt{117} \approx 10{,}817$}{}
\lz{c)}{$100$ cm $= 1$ m}{$0{,}6$ m $= 60$ cm; \ $60^2 + 80^2 = 10\,000$ \ $\Rightarrow$ \ $\sqrt{10\,000} = 100$}{}
\end{pfloesung}
\begin{pfloesung}{}
\lz{7.}{ja}{Leiter bis zum Fenster: $\sqrt{0{,}7^2 + 2{,}4^2} = \sqrt{6{,}25} = 2{,}5$ m $< 2{,}6$ m}{}
\end{pfloesung}
\begin{pfloesung}{}
\lz{8.}{$40$ m}{quer: $\sqrt{80^2 + 60^2} = 100$ m; \ $140 - 100 = 40$}{}
\end{pfloesung}
\begin{pfloesung}{}
\lz{9.}{$\sqrt{8\,425} \approx 91{,}8$ cm}{$80^2 + 45^2 = 8\,425$ \ $\Rightarrow$ \ $\approx 91{,}79$}{}
\end{pfloesung}
\begin{pfloesung}{}
\lz{10.}{$\overline{AB} = \sqrt{1\,300} \approx 36{,}1$ m}{$\overline{AB}^2 = 12^2 + 34^2 = 1\,300$ \ $\Rightarrow$ \ $\approx 36{,}06$}{}
\end{pfloesung}
\begin{pfloesung}{}
\lz{11.}{$x = \sqrt{29\,156} \approx 170{,}8$ cm}{$x^2 = 170^2 + 16^2 = 29\,156$ \ $\Rightarrow$ \ $\approx 170{,}75$}{}
\end{pfloesung}
\end{document}
